\section{Evaluation protocol}
\label{sec:protocol}

\subsection{Standing rules}

Five rules applied to every model in the project.
\begin{description}[leftmargin=1.2em,style=nextline]
  \item[Vantage.] Anything used to predict a game in season $V$ is computed
  from seasons before $V$ and from games of season $V$ played strictly before
  it: scalers, calibration maps, embeddings, priors and coordinate
  normalisers included. The dressed lineup and the announced starting goalie
  are admissible, because both are public about an hour before the game.
  \item[Market firewall.] No odds or market prices in any feature, target,
  calibration or selection step.
  \item[External models.] Quantities fitted by third parties on full-sample
  data are excluded; only raw recorded fields are admitted.
  \item[Era conditioning.] Era enters through walk-forward covariates (league
  scoring to date, overtime format, shootout and loser-point flags, a scaled
  season index), never through season indicators or full-sample
  normalisation.
  \item[Determinism.] Seeds are fixed and recorded. Neural training runs on
  Apple MPS, which is not bit-deterministic, so each trained checkpoint is
  identified by its SHA-256, every reported number comes from CPU inference
  over a saved checkpoint, and committed per-game predictions must re-derive
  within $10^{-6}$.
\end{description}

\subsection{Windows}

Seasons were assigned to roles before any model was fitted (Table~\ref{tab:data},
Figure~\ref{fig:windows}). For \neurhl{} 1.0 the windows are development
(2009--2011), tune (2012--2017) and confirmation (2018--2026, touched once per
layer), followed by the live 2027 season. Player-level layers use their own
split: development on 2011--2012, tuning on 2014--2017, a declared evaluation
window of 2022--2026 and a one-shot confirmation on 2018--2020. \neurhlG{}, built
after the 1.0 confirmation had spent 2018--2026 at game level, uses a fresh
split of the same history: iteration on 2012 and 2014--2018 (7,421 games),
gates on 2019, 2020 and 2022--2024 (6,289 games), and a sealed test on
2025--2026 (2,624 games). Those seasons had already been scored at game level
by \neurhlH{}, so the \neurhlG{} seal protects against adaptive \neurhlG{} choices; it is
procedurally sealed, not unseen, and we treat the 2026--27 live season as the
only clean test of \neurhlG{}.

The windows are enforced in code. A scoring guard refuses to score an
anomalous season or to cross a window boundary. For \neurhlG{}, a single guarded
loader is the only route by which model, training and evaluation code may
read the data; it drops the sealed seasons unless an unseal function has been
called, and that function refuses every caller except the seal script. A
syntax-tree test fails any \neurhlG{} module that reads the tensor directory
directly, and the 37 sealed input files are pinned by hash.

\begin{figure}[t]
\centering
\includegraphics[width=\textwidth]{fig_windows}
\caption{Season roles in the two protocols. \neurhl{} 1.0 spent its
confirmation window, 2018--2026, once per layer. \neurhlG{} reuses the same history
under a new split whose sealed seasons, 2025--2026, remain unspent. Lockout
(2013) and COVID (2021) seasons train but are never scored.}
\label{fig:windows}
\end{figure}

\subsection{Commit before results}

Every plan, amendment and candidate declaration was committed before the
numbers that tested it existed, and every freeze before the confirmatory run it
governs. Amendments are appended, never edited. The repository was not public
before 26 September 2026, so for the plans written before then the ordering
rests on the commit times recorded in its history, which a reader has to take
on trust. From publication on, each push to the public repository is dated
independently of the author: the \neurhlG{} plan and its candidate declaration
were both on the public branch before any \neurhlG{} gate number existed, and
every live forecast must be there before its game starts. Each plan fixes
its gates with numbers, its search budget and the consequence of every
outcome. Searches are logged row by row in ledgers with hard caps: 40
configurations for the first neural game model, 12 for the game-level work of
the second and third plans, 10, 4, 4 and 4 for the player-game, player-season,
goalie and expected-goals layers, and for \neurhlG{} at most 14 full runs on the
iteration window and at most two on the gate window, the second only as one
declared change after a failure. Exceeding a cap voids the version.

Confirmatory tests are one-shot. Each has its own script that refuses to run
before its freeze commit is on the public branch, refuses to run twice, and
writes its result in a single step. A pre-gate on non-confirmatory seasons
decides whether a confirmatory window is spent at all, so a weak candidate
cannot consume it.

\subsection{Causality audits}

Gates that score a model on the inputs it was trained with cannot detect an
illegitimate input: a leaked feature improves in-sample realism along with
everything else (Section~\ref{sec:failures}). Every data builder therefore
passes a causality audit before any gated number is reported from its output.
The audit corrupts every row after a cut date, rebuilds the features, and
requires every value before the cut to be bit-identical. For the event
simulator the same test runs at the level of individual model inputs, together
with a target-identity check (no input may agree with any target above 98\%)
and an empirical check that exempted inputs carry no lift on the next event.
The \neurhlG{} state builders are audited the same way, with cuts in mid-season
2016 and at the boundary between the last gate season and the first sealed
one; all 100 skater-state, 23 goalie-state and 53 team-state features are
bit-identical before each cut.

\subsection{Statistics}

The primary metric for win probabilities is per-game log loss, compared
between models as a mean paired difference over the same games. We report a
normal-approximation test on the per-game differences, a season-clustered $t$
test, a wild cluster bootstrap on seasons that enumerates all $2^S$
Rademacher sign patterns exactly, a moving-block bootstrap over date-ordered
games (blocks of 100 games, 9,999 draws, seed 711), a sign test across
seasons, and the Murphy decomposition of the Brier score into reliability,
resolution and uncertainty. A log-loss gain with no gain in resolution is
reported as recalibration, not information. Player-level differences use
standard errors clustered two ways, by player and by game
\citep{cameron2011robust}, and families of tests are Holm-corrected. Count
predictions are assessed with Poisson deviance, the continuous ranked
probability score and randomised probability integral transforms.

Power is stated before each test. On the tune window, one standard error of a
paired log-loss difference against Elo is 0.00268 nats. A gate of the form
``beat Elo by 0.002'' is therefore a 0.75-standard-error bar, which a model
with no skill clears by chance 23\% of the time. The first plan's second
amendment made a significance test the headline criterion for exactly this
reason, and every later plan kept it.
